\chapter{Enriched Category Theory}
\label{ch:enriched}

\section{Motivation}

In an ordinary locally small category, hom-sets $\cat{C}(a,b)$ are sets. But in many mathematical contexts, hom-sets carry additional structure: they are abelian groups (additive categories), vector spaces (linear categories), topological spaces (topological categories), or simplicial sets (simplicial categories). Enriched category theory provides a unified framework for all these situations.

The key idea: replace $\Set$ as the ``ground category'' for hom-sets with a more structured monoidal category $\cat{V}$. Instead of hom-sets, we have hom-\emph{objects} in $\cat{V}$.

\section{Monoidal Categories}

\begin{definition}[Monoidal category]
A \textbf{monoidal category} is a category $\cat{V}$ equipped with:
\begin{itemize}
  \item A \textbf{tensor product} bifunctor $\otimes: \cat{V} \times \cat{V} \to \cat{V}$.
  \item A \textbf{unit object} $I \in \cat{V}$.
  \item Natural isomorphisms:
  \begin{itemize}
    \item \textbf{Associator}: $\alpha_{A,B,C}: (A \otimes B) \otimes C \xrightarrow{\sim} A \otimes (B \otimes C)$.
    \item \textbf{Left unitor}: $\lambda_A: I \otimes A \xrightarrow{\sim} A$.
    \item \textbf{Right unitor}: $\rho_A: A \otimes I \xrightarrow{\sim} A$.
  \end{itemize}
\end{itemize}
satisfying the pentagon axiom (coherence for the associator) and the triangle axiom (coherence for the unitors).

$\cat{V}$ is \textbf{symmetric monoidal} if it additionally has a natural isomorphism $\sigma_{A,B}: A \otimes B \xrightarrow{\sim} B \otimes A$ satisfying the hexagon axioms and $\sigma_{B,A} \circ \sigma_{A,B} = \id$.
\end{definition}

\begin{theorem}[Mac Lane's Coherence Theorem]
In any monoidal category, all diagrams built from associators and unitors commute. Thus we may treat monoidal categories as if they were strictly associative and unital (up to coherent isomorphism).
\end{theorem}

\begin{example}[Standard monoidal categories for enrichment]
\leavevmode
\begin{itemize}
  \item $(\Set, \times, \{*\})$: ordinary categories.
  \item $(\Ab, \otimes_\mathbb{Z}, \mathbb{Z})$: preadditive ($\Ab$-enriched) categories.
  \item $(\Vect_k, \otimes_k, k)$: $k$-linear categories.
  \item $(\mathbf{Cat}, \times, \mathbf{1})$: strict 2-categories.
  \item $(\mathbf{sSet}, \times, \Delta^0)$: simplicially enriched categories.
  \item $([0,\infty], +, 0)$ (Lawvere): extended metric spaces.
  \item $(\{T, F\}, \wedge, T)$ (the Boolean poset): preorders.
\end{itemize}
\end{example}

\section{$\cat{V}$-Enriched Categories}

\begin{definition}[$\cat{V}$-enriched category]
Let $(\cat{V}, \otimes, I)$ be a monoidal category. A \textbf{$\cat{V}$-enriched category} (or \textbf{$\cat{V}$-category}) $\cat{C}$ consists of:
\begin{itemize}
  \item A collection $\ob(\cat{C})$ of objects.
  \item For each pair $a, b \in \ob(\cat{C})$, an object $\cat{C}(a,b) \in \cat{V}$ (the hom-object).
  \item For each triple $a,b,c$, a \textbf{composition morphism} in $\cat{V}$:
  \[
    \circ_{a,b,c}: \cat{C}(b,c) \otimes \cat{C}(a,b) \to \cat{C}(a,c).
  \]
  \item For each object $a$, an \textbf{identity morphism} in $\cat{V}$:
  \[
    j_a: I \to \cat{C}(a,a).
  \]
\end{itemize}
subject to the usual associativity and unit axioms (expressed as commutative diagrams in $\cat{V}$).
\end{definition}

\begin{remark}
An ordinary (locally small) category is precisely a $\Set$-enriched category. The hom-object is the hom-set; composition and identities are functions.
\end{remark}

\begin{example}[$\Ab$-enriched categories (preadditive)]
An $\Ab$-enriched category has abelian groups as hom-objects and bilinear composition. Examples: $\Ab$ itself (as a $\Ab$-enriched category), $\Vect_k$, $\mathbf{Mod}_R$. Any additive category is an $\Ab$-enriched category with a zero object, finite products, and finite coproducts (which coincide).
\end{example}

\begin{example}[Lawvere metric spaces]
A $([0,\infty], +, 0)$-enriched category is precisely a (generalized) metric space in Lawvere's sense: objects are points; $\cat{C}(a,b) \in [0,\infty]$ is the distance; the composition axiom is the triangle inequality $d(a,c) \leq d(b,c) + d(a,b)$; the identity axiom is $d(a,a) \leq 0$, i.e., $d(a,a) = 0$. Non-symmetry and non-Hausdorff metric spaces are naturally captured. Enriched functors are contractive maps; enriched natural transformations are non-expansive morphisms.
\end{example}

\begin{example}[2-categories as $\Cat$-enriched categories]
A $\Cat$-enriched category is a \textbf{2-category}: its hom-objects are categories (whose objects are the 1-cells and whose morphisms are the 2-cells), and composition is a functor (giving horizontal composition of 2-cells).
\end{example}

\begin{example}[Simplicial categories]
A $\mathbf{sSet}$-enriched category (simplicial category) has simplicial sets as hom-objects. These arise in homotopy theory and higher category theory: the mapping spaces encode homotopy types of spaces of morphisms.
\end{example}

\section{$\cat{V}$-Functors and $\cat{V}$-Natural Transformations}

\begin{definition}[$\cat{V}$-functor]
Let $\cat{C}$ and $\cat{D}$ be $\cat{V}$-categories. A \textbf{$\cat{V}$-functor} $F: \cat{C} \to \cat{D}$ consists of:
\begin{itemize}
  \item A function $F: \ob(\cat{C}) \to \ob(\cat{D})$.
  \item For each $a, b \in \cat{C}$, a morphism in $\cat{V}$: $F_{a,b}: \cat{C}(a,b) \to \cat{D}(Fa,Fb)$.
\end{itemize}
preserving composition and identities (as commutative diagrams in $\cat{V}$).
\end{definition}

\begin{definition}[$\cat{V}$-natural transformation]
A \textbf{$\cat{V}$-natural transformation} $\alpha: F \Rightarrow G$ (between $\cat{V}$-functors $F, G: \cat{C} \to \cat{D}$) consists of morphisms $\alpha_a: I \to \cat{D}(Fa, Ga)$ in $\cat{V}$ (``elements'' of the hom-objects) satisfying the enriched naturality condition:
\[
  \begin{tikzcd}
    \cat{C}(a,b) \arrow[r, "F_{a,b}"] \arrow[d, "G_{a,b}"'] & \cat{D}(Fa,Fb) \arrow[d, "(\alpha_b)_*"] \\
    \cat{D}(Ga,Gb) \arrow[r, "(\alpha_a)^*"'] & \cat{D}(Fa,Gb)
  \end{tikzcd}
\]
(where $(\alpha_b)_*$ and $(\alpha_a)^*$ are composition with $\alpha_b$ and $\alpha_a$, using the composition morphisms of $\cat{D}$).
\end{definition}

\begin{remark}
When $\cat{V} = \Set$, a $\cat{V}$-natural transformation reduces to an ordinary natural transformation (since $\alpha_a: \{*\} \to \cat{D}(Fa, Ga)$ picks out a single element of the hom-set).
\end{remark}

$\cat{V}$-categories, $\cat{V}$-functors, and $\cat{V}$-natural transformations form a 2-category $\cat{V}$-$\Cat$.

\section{The Enriched Yoneda Lemma}

\begin{theorem}[Enriched Yoneda Lemma]
Let $\cat{C}$ be a $\cat{V}$-category and $F: \cat{C} \to \cat{V}$ a $\cat{V}$-functor (where $\cat{V}$ is enriched over itself via its closed structure, when $\cat{V}$ is closed monoidal). Then:
\[
  [\cat{C}, \cat{V}](\cat{C}(a,-), F) \;\cong\; Fa
\]
as objects of $\cat{V}$, naturally in $a \in \cat{C}$ and $F \in [\cat{C},\cat{V}]$.
\end{theorem}

The enriched Yoneda lemma requires $\cat{V}$ to be \textbf{closed monoidal}: for each $A \in \cat{V}$, the functor $A \otimes -$ has a right adjoint $[A, -]$ (the \textbf{internal hom}). Key examples:
\begin{itemize}
  \item $(\Set, \times)$: internal hom is $[A,B] = B^A$ (function set).
  \item $(\Ab, \otimes_\mathbb{Z})$: internal hom is $\Hom_\mathbb{Z}(A,B)$ (abelian group of homomorphisms).
  \item $(\Vect_k, \otimes_k)$: internal hom is $\Hom_k(V, W)$.
  \item $(\mathbf{sSet}, \times)$: internal hom is the simplicial mapping space $\mathrm{Map}(X,Y)$.
\end{itemize}

\section{Enriched Limits: Weighted Limits}

Ordinary limits are controlled by shape categories $\cat{J}$. In the enriched setting, the shape data includes a \textbf{weight} functor.

\begin{definition}[Weighted limit]
Let $\cat{C}$ be a $\cat{V}$-category, $D: \cat{J} \to \cat{C}$ a $\cat{V}$-functor, and $W: \cat{J} \to \cat{V}$ a weight functor. The \textbf{$W$-weighted limit} of $D$ is an object $\{W, D\} \in \cat{C}$ representing the $\cat{V}$-functor:
\[
  c \;\mapsto\; [\cat{J}, \cat{V}](W, \cat{C}(c, D-)): \cat{C}\op \to \cat{V}.
\]
That is, $\cat{C}(c, \{W,D\}) \cong [\cat{J},\cat{V}](W, \cat{C}(c, D-))$ in $\cat{V}$, naturally in $c$.
\end{definition}

\begin{remark}
Ordinary (conical) limits correspond to the constant weight $W = \Delta_I$ (the constant functor at the monoidal unit $I$). Weighted limits strictly generalize ordinary limits and are essential in enriched category theory, where conical limits are often insufficient.
\end{remark}

\begin{example}[Ends as weighted limits]
The \textbf{end} $\int_c T(c,c)$ of a $\cat{V}$-functor $T: \cat{C}\op \otimes \cat{C} \to \cat{V}$ (where $\cat{C}\op \otimes \cat{C}$ is the tensor product of $\cat{V}$-categories) is the limit of $T$ weighted by the hom-functor $\cat{C}(-,-): \cat{C}\op \otimes \cat{C} \to \cat{V}$.

The $\cat{V}$-valued natural transformations $\Nat(F,G)$ are an end:
\[
  \Nat(F,G) = \int_c \cat{D}(Fc, Gc).
\]
\end{example}

\section{Change of Base}

\begin{definition}[Change of base]
A \textbf{lax monoidal functor} $\Phi: \cat{V} \to \cat{W}$ (with coherence maps $\phi_{A,B}: \Phi A \otimes_{\cat{W}} \Phi B \to \Phi(A \otimes_{\cat{V}} B)$ and $\phi_I: I_{\cat{W}} \to \Phi(I_{\cat{V}})$) induces a \textbf{change of base} 2-functor $\Phi_*: \cat{V}\text{-}\Cat \to \cat{W}\text{-}\Cat$: apply $\Phi$ to all hom-objects and compose the structural maps.
\end{definition}

\begin{example}
The ``set of connected components'' functor $\pi_0: \mathbf{sSet} \to \Set$ is lax monoidal. Change of base along $\pi_0$ sends a simplicial category to its underlying ordinary category (obtained by taking connected components of each hom-space).

The free abelian group functor $\mathbb{Z}[-]: \Set \to \Ab$ is lax monoidal. Change of base sends ordinary categories to $\Ab$-enriched categories (the ``abelianization'' of a category).
\end{example}

\section{Exercises}

\begin{exercise}
Verify that an $[0,\infty]$-enriched category is precisely a generalized metric space (in the sense of Lawvere): a set $X$ with a function $d: X \times X \to [0,\infty]$ satisfying $d(x,x) = 0$ and the triangle inequality.
\end{exercise}

\begin{exercise}
Show that in any $\Ab$-enriched category, for any two objects $a, b$, the hom-group $\cat{C}(a,b)$ carries an abelian group structure compatible with composition (i.e., composition is bilinear).
\end{exercise}

\begin{exercise}
Define the $\cat{V}$-enriched functor category $[\cat{C},\cat{D}]$ for $\cat{V}$-categories $\cat{C}$ and $\cat{D}$ (assuming $\cat{V}$ is closed monoidal). What are its hom-objects?
\end{exercise}

\begin{exercise}[$\star$]
Let $\cat{C}$ be a $\cat{V}$-category with $\cat{V} = \Ab$. Define the notion of a $\cat{V}$-functor $F: \cat{C} \to \Ab$ and show that $\Ab$-natural transformations between such functors form an abelian group.
\end{exercise}

\begin{exercise}[$\star$]
Prove that for a small $\Ab$-enriched category $\cat{C}$, the functor category $[\cat{C}\op, \Ab]$ is an abelian category (assuming familiarity with abelian categories). Identify the representable functors with projective objects.
\end{exercise}

\begin{exercise}[$\star\star$]
Formulate and prove the enriched Yoneda lemma for $\cat{V}$-functors $F: \cat{C} \to \cat{V}$ where $\cat{V}$ is a symmetric closed monoidal category. Show how it reduces to the ordinary Yoneda lemma when $\cat{V} = \Set$.
\end{exercise}
